\newcommand{\przyrzadopis}[1]{%
\par\medskip\noindent
\begin{minipage}[c]{0.44\linewidth}
\centering\includegraphics[width=\linewidth,height=104mm,keepaspectratio]{#1}
\end{minipage}\hfill
\begin{minipage}[c]{0.51\linewidth}
\raggedright
\captionof{figure}{Przyrząd do statycznego pomiaru modułu Younga.}
\medskip
\textbf{A — górny uchwyt}\par
Mocuje górny koniec badanego drutu.
\par\smallskip\textbf{B — dolne mocowanie}\par
Łączy drut z układem obciążenia.
\par\smallskip\textbf{C — dźwignia}\par
Przenosi przemieszczenie drutu na czujnik.
\par\smallskip\textbf{D — czujnik mikrometryczny}\par
Wskazuje przemieszczenie dwukrotnie większe od wydłużenia drutu.
\par\smallskip Szalka z odważnikami obciąża dolne mocowanie.
\medskip\par{\footnotesize Rysunek: instrukcja \texttt{11\_wykon.pdf}, s.~2.}
\end{minipage}\par\medskip}
\newcommand{\teoria}{%
W zakresie liniowej sprężystości prawo Hooke'a wiąże naprężenie
$\sigma$ z odkształceniem względnym $\varepsilon$:
\[
\boxed{\sigma=E\varepsilon,\quad\sigma=\frac{F}{S},\quad
\varepsilon=\frac{\Delta l}{l},\quad S=\frac{\pi d^2}{4}.}
\]
Moduł Younga $E$ ma jednostkę Pa. Zależność wydłużenia od siły
w przyjętym zakresie opisujemy prostą ze swobodnym wyrazem wolnym:
\[
\boxed{\Delta l=aF+b,\qquad E=\frac{4l}{\pi d^2a},\qquad F=mg.}
\]
Wzór na $E$ wymaga jednostek SI: $l,d$ w m oraz $a$ w m/N.
Dla dopasowania względem masy $\Delta l=cm+b$ zamiast siły
należy użyć $E=4lg/(\pi d^2c)$; nie wolno zamieniać współczynników $a$ i $c$.
}
\newcommand{\dzwignia}{%
Czujnik jest sprzężony z drutem przez dźwignię. Odległość drutu
od osi obrotu jest dwa razy mniejsza niż odległość punktu styku czujnika,
zatem jego wskazanie jest dwa razy większe od wydłużenia drutu.
Po prawidłowym wyzerowaniu:
\[
\boxed{\Delta l_{\uparrow}=\frac{z_{\uparrow}}2,\quad
\Delta l_{\downarrow}=\frac{z_{\downarrow}}2,\quad
\overline{\Delta l}=\frac{z_{\uparrow}+z_{\downarrow}}4.}
\]
Jeśli zero nie jest równe 0, przed użyciem wzoru odejmij odpowiedni
odczyt zerowy. Nie dziel wskazań przez 2 ponownie w regresji.
}
\newcommand{\tabelapomiarow}{%
\begingroup\small\renewcommand{\arraystretch}{1}
\begin{tabularx}{\linewidth}{|c|C|C|C|C|Y|}\hline
\naglowek\textbf{Nr} & \textbf{Masa całkowita} [kg] & \textbf{Czujnik $z_\uparrow$} [mm] &
\textbf{Czujnik $z_\downarrow$} [mm] & \textbf{Zero / uwagi} & \textbf{Uwagi do punktu}\\\hline
\csname @@input\endcsname{tabela_pusta.tex}
\end{tabularx}\endgroup}
\newcommand{\tabelaopracowania}{%
\begingroup\small\renewcommand{\arraystretch}{1}
\begin{tabularx}{\linewidth}{|c|C|C|C|C|C|}\hline
\naglowek\textbf{Nr} & $m$ [kg] & $F=mg$ [N] &
$z_\uparrow$ [mm] & $z_\downarrow$ [mm] & $\overline{\Delta l}$ [mm]\\\hline
\csname @@input\endcsname{tabela_pusta.tex}
\end{tabularx}\endgroup}
\newcommand{\wymiary}[1]{%
\textbf{#1}\quad Materiał: \pole\quad Długość $l=$ \pole\ m.
\par\smallskip
\begin{tabularx}{\linewidth}{|Y|C|C|C|}\hline
\naglowek\textbf{Średnica w trzech miejscach} & $d_1$ [mm] & $d_2$ [mm] & $d_3$ [mm]\\\hline
Pomiar przy obciążeniu 2 kg & & &\rule{0pt}{9mm}\\\hline
\end{tabularx}\par\smallskip}
